\title{From Data to Cure: A Comprehensive Exploration of Multi-omics Data Analysis for Targeted Therapies}
\begin{document}
\maketitle

\begin{abstract}
In the dynamic landscape of targeted therapeutics, drug discovery has pivoted towards understanding underlying disease mechanisms, placing a strong emphasis on molecular perturbations and target identification. This paradigm shift, crucial for drug discovery, is underpinned by big data, a transformative force in the current era. Omics data, characterized by its heterogeneity and enormity, has ushered biological and biomedical research into the big data domain. Acknowledging the significance of integrating diverse omics data strata, known as multi-omics studies, researchers delve into the intricate inter- relationships among various omics layers. This review navigates the expansive omics landscape, showcasing tailored assays for each molecular layer through genomes to metabolomes. The sheer volume of data generated necessitates sophisticated informatics techniques, with machine-learning (ML) algorithms emerging as robust tools. These datasets not only refine disease classification but also enhance diagnostics and foster the development of targeted therapeutic strategies. Through the integration of high-throughput data, the review focuses on targeting and modeling multiple disease-regulated networks, validating interactions with multiple targets, and enhancing therapeutic potential using network pharmacology approaches. Ultimately, this exploration aims to illuminate the transformative impact of multi-omics in the big data era, shaping the future of biological research.
\end{abstract}

\section{Recovered Text}
Vol.:(0123456789)
Molecular Biotechnology (2025) 67:1269–1289
https://doi.org/10.1007/s12033-024-01133-6
REVIEW PAPER
From Data to Cure: A Comprehensive Exploration of Multi‑omics Data
Analysis for Targeted Therapies
Arnab Mukherjee1 · Suzanna Abraham1 · Akshita Singh1 · S. Balaji1 · K. S. Mukunthan1
Received: 27 December 2023 / Accepted: 27 February 2024 / Published online: 2 April 2024
© The Author(s) 2024
Abstract
In the dynamic landscape of targeted therapeutics, drug discovery has pivoted towards understanding underlying disease
mechanisms, placing a strong emphasis on molecular perturbations and target identification. This paradigm shift, crucial
for drug discovery, is underpinned by big data, a transformative force in the current era. Omics data, characterized by its
heterogeneity and enormity, has ushered biological and biomedical research into the big data domain. Acknowledging the
significance of integrating diverse omics data strata, known as multi-omics studies, researchers delve into the intricate inter-
relationships among various omics layers. This review navigates the expansive omics landscape, showcasing tailored assays
for each molecular layer through genomes to metabolomes. The sheer volume of data generated necessitates sophisticated
informatics techniques, with machine-learning (ML) algorithms emerging as robust tools. These datasets not only refine
disease classification but also enhance diagnostics and foster the development of targeted therapeutic strategies. Through
the integration of high-throughput data, the review focuses on targeting and modeling multiple disease-regulated networks,
validating interactions with multiple targets, and enhancing therapeutic potential using network pharmacology approaches.
Ultimately, this exploration aims to illuminate the transformative impact of multi-omics in the big data era, shaping the
future of biological research.
Keywords  Multi-omics · Big data · Machine learning · Targeted therapeutics · Network pharmacology
Introduction
In the era of targeted therapeutics, drug discovery
approaches emphasize the underlying disease mechanisms,
accompanied by target identification and lead discovery.
This targeted therapeutic system focuses on molecular
perturbations, has become critical in the drug discovery
process [1], and is entirely reliant on big data to transform
the current understanding and accessible data into valuable
information employed to enhance clinical outcomes [2].
Biological and biomedical research and its applications
have infiltrated a big data era due to the heterogeneity and
enormity of omics data [3]. There is a growing acknowl-
edgment among researchers regarding the significance of
adeptly integrating various strata of omics data, referred
to as multi-omic studies. The modeling and exploration of
the intricate interrelationships among diverse omic layers
have the potential to unveil crucial functional and clinical
insights. The transformative omics revolution observed
in biological research after the inception of genomic
sequencing has engendered an extensive corpus of data,
concurrently fostering technologies that facilitate the
cost-effective and streamlined quantification of biologi-
cal molecules on a large scale. Presently, omics technolo-
gies have undergone substantial expansion to encompass
more unbounded methodologies, such as assays reliant on
next-generation sequencing (NGS) and mass spectrometry.
Tailored assays now exist for each stratum of molecular
 *	 K. S. Mukunthan
mukunthan.ks@manipal.edu; mukunthanselvam@gmail.com
Arnab Mukherjee
arnabbiotech.gen@gmail.com
Suzanna Abraham
suzanna16abraham@yahoo.com
Akshita Singh
akshitasingh9156@gmail.com
S. Balaji
s.balaji@manipal.edu
1
Department of Biotechnology, Manipal Institute
of Technology, Manipal Academy of Higher Education,
Manipal, India
1270
Molecular Biotechnology (2025) 67:1269–1289
activity, spanning genomes to metabolomes. Researchers
investigate the field of genomics through techniques like
whole-genome or whole-exome sequencing, explore tran-
scriptomics using RNA-seq [4], delve into epigenomics via
methodologies such as bisulfite sequencing (BS-seq) [5],
ChIP-seq for histone modifications [6], and ATAC-Seq for
open chromatin [7]. The three-dimensional conformation
of the genome can be elucidated through techniques like
Hi-C or chromatin interaction analysis with paired-end tag
(ChIA-PET) [8, 9].
Additionally, researchers may investigate proteomics
and delve into metabolomics, predominantly employing
mass spectrometry [10]. These methodologies have trans-
formed biomedical investigations by furnishing a more
thorough understanding of the studied biological system
and the molecular intricacies inherent in the progression
of diseases. The prodigious volume of data generated in
biomedicine necessitates the employment of sophisticated
informatics techniques to glean novel insights, advance our
understanding of diseases, enhance diagnostic capabilities,
and formulate individualized therapeutic strategies [11].
Within this framework, machine-learning (ML) algorithms
have emerged as among the most auspicious methods in
the field. Omics data have the potential to refine classifi-
cation beyond a simplistic dichotomy of healthy versus
diseased, offering substantial clinical benefits. These
methodologies can enhance patient treatment by align-
ing with the distinctive biology of their specific ailment,
categorizing patients into subtypes, or positioning them
along a spectrum of disease manifestation. These datasets
also contribute significantly to an enhanced comprehen-
sion of the pathogenic mechanisms underlying diseases
and biomarkers. We briefly reviewed how the accessibility
of these omics data have transformed biology and aided in
developing systems biology to comprehend the biological
phenomena [12] and provide a platform for integrating
these high-throughput data, which focuses on targeting and
modeling multiple disease-regulated networks to screen
leads and validating their interactions with multiple targets
for enhanced therapeutic potential [13].
This review aims to explore the transformative impact
of omics data in the big data era of biological and bio-
medical research. Focusing on multi-omic studies, it
seeks to model and understand the intricate interrelation-
ships among diverse omic layers to unveil functional and
clinical insights. The expansion of omics technologies
has provided tailored assays for each molecular stratum,
revolutionizing biomedical investigations. Leveraging ML
algorithms, the study aims to refine disease classification,
enhance diagnostics, and develop personalized therapeutic
strategies, ultimately contributing to an improved under-
standing of pathogenic mechanisms and biomarkers.
Multi‑omics Data
Genomics
Genomic investigations have constituted a principal
methodology in elucidating the etiology of diseases and
delineating potential targets for treating various complex
diseases. The statistical analyses conducted in these stud-
ies frequently encounter challenges related to multiplic-
ity, faint signals, and the inherent interdependence among
genetic markers [14]. ML algorithms can be implemented
to address these challenges. These algorithms uncover sub-
tle patterns and relationships by improving the sensitiv-
ity and specificity of the analysis that may not be evident
through conventional statistical approaches. Yu et al. illus-
trated the effectiveness of an integrative co-localization
(INCO) algorithm designed for the seamless integration
of single-nucleotide variants (SNVs) and copy number
variations (CNVs). This algorithm facilitated the synthe-
sis of the genetic variations, yielding a more precise and
refined genetic region. The refined region enhanced the
accuracy in identifying causal variants associated with
the studied biological phenomena [15]. This allows the
identification of specific genetic mutations or alterations
associated with diseases. Also, Liu et al. demonstrated a
comprehensive analysis of these genetic variations, which
resulted in the identification of disease-causing variants
in 10 of the 16 investigated rare diseases. Notably, the
analysis revealed new potentially pathogenic variants for
two disorders. For the first time, clinical whole genome
sequencing (WGS) successfully identified a causative
simple sequence repeat (SSR) variation associated with
Machado–Joseph disease, highlighting the power of clini-
cal WGS in providing molecular-level diagnostic clarity
for rare diseases [16]. By unraveling the intricate genomic
landscape, targeted therapies can be precisely tailored to
address the underlying genetic anomalies, paving the way
for personalized treatment strategies that consider indi-
vidual genetic variations.
Transcriptomics
Transcriptomics pertains to the assessment of the entire
repertoire of genes expressed as transcripts or mRNAs
and non-coding RNAs [17]. Alterations in the gene
expression profile often occur under diseased conditions.
Transcriptomic methodologies, such as microarrays and
high-throughput sequencing, facilitate the systematic
monitoring of the entire transcriptome [18]. This com-
prehensive approach allows for the acquisition of a global
cellular signature or fingerprint, offering insights into
1271
Molecular Biotechnology (2025) 67:1269–1289
the dynamic changes in gene expression associated with
various biological states [19]. Moreover, it establishes a
critical bridge between genomics and proteomics by eluci-
dating the intricate connection between the transcriptome
and the subsequent protein expression patterns [20]. In
targeted therapies, understanding transcriptomic changes
is crucial for identifying key genes and pathways involved
in disease progression. By employing these techniques,
researchers unveiled the dysregulated genes and designed
interventions that modulate gene expression levels, ulti-
mately contributing to the development of more effective
and tailored therapeutic approaches [21–23].
Epigenomics
The dysregulation of epigenetic processes is a pivotal factor
in the onset and advancement of human diseases. Due to
the dynamic nature of the intricately regulated epigenetic
marks and mechanisms, these modifications can serve as
discernible biomarkers [24]. Differential DNA methylation
can lead to a spectrum of disorders, encompassing inflam-
matory conditions, precancerous lesions, and malignancies
[25–27]. Moreover, Kikutake and Yahara, in 2016, con-
ducted a genome-wide study that illustrated the advantage
of ChIP-Seq and RNA-Seq over microarrays in encompass-
ing histone modification regions not addressed by microar-
rays. The study delved into the exploration of associations
between histone modifications and the occurrence of aber-
rant gene expression during the progression of the disease
[28]. Comprehending epigenetic modifications is crucial for
identifying reversible alterations that influence gene expres-
sion in targeted therapeutics. A more sophisticated approach
to therapy can be achieved by developing medicines that
selectively affect the activity of genes linked to disease pro-
gression by focusing on these epigenetic alterations [29].
Proteomics
Aberrant regulation of protein function plays a critical role
in disease pathogenesis, highlighting the imperative goal of
biomedical research to comprehend the perturbation of the
proteome in disease progression. While transcriptome data,
specifically mRNA abundance, cannot infer protein abun-
dance accurately, direct assessments of protein function
become essential [30]. Although conventional methodolo-
gies often concentrate on individual proteins or a limited set,
recent progress in sample separation and mass spectrometry
technologies allows for the holistic consideration of a com-
plex biological system as an integrated unit [31]. The swift
progress in proteomics experimental techniques has spurred
the development of diverse downstream bioinformatics
analyses. These analyses contribute significantly to eluci-
dating the intricate relationships between molecular-level
protein regulatory mechanisms and phenotypic manifesta-
tions, particularly in the context of disease initiation and
progression [32]. Ohlsson et al. analyzed differential protein
expression, utilizing statistical methods to identify analytes
specific to autoimmune diseases that perturb immunoregu-
latory responses [33]. By identifying dysregulated proteins,
researchers can design therapies that specifically target these
molecular players, addressing key components of the disease
mechanisms. This nuanced understanding of the proteomic
landscape enhances the precision and efficacy of targeted
therapeutic interventions.
Metabolomics
Metabolic profiles offer a detailed understanding of physi-
ological states and are highly susceptible to genetic and envi-
ronmental perturbations. Fluctuations in metabolic profiles
can offer insights into the mechanisms underlying patho-
logical conditions, thereby serving as potential biomarkers
for the diagnosis and evaluation of the risk associated with
disease onset [34]. Large-scale metabolomics data sources
are abundant as high-throughput technologies continue to
evolve. For instance, meticulous statistical and bioinfor-
matic analysis of intricate metabolomics data are crucial for
achieving accurate and significant findings that are applied
in real-world clinical settings [35, 36]. Metabolites serve as
biomarkers and contribute to an enhanced comprehension
of the pathophysiology underlying various diseases [37–40].
In targeted therapies, metabolomics helps identify disease-
associated metabolic signatures, shedding light on altered
biochemical pathways [38, 41]. This knowledge guides the
development of therapeutic agents that target-specific meta-
bolic processes, addressing the unique metabolic needs of
diseased cells and contributing to more efficient and accurate
therapeutic strategies.
Exploitation of Omics Data
The high-dimensional structure of omics data raises several
barriers to acquiring information. Data processing, normali-
zation, integration, and analysis of these high-dimensional
data have gained attention among researchers through sev-
eral computational techniques such as ML, meta-heuristic,
and statistical approaches [42]. The exponential increase
in the volume of data generated through high-throughput
sequencing and related methodologies necessitates the appli-
cation of statistical models capable of extracting precise
and interpretable predictions from the wealth of biological
information [11]. ML models exhibit enhanced performance
when trained on large datasets, making them well-suited for
the intricate integration of multi-omics data in bioinformat-
ics applications [43]. The speed, accuracy, interpretability,
1272
Molecular Biotechnology (2025) 67:1269–1289
computing cost, complexity, and sample sizes were prag-
matically scored in the range of 1–4, with low to very high as
a threshold for supervised ML algorithms by Reel et al. [44],
as demonstrated in Fig. 1. These ML models have accel-
erated the drug development process by identifying novel
biomarkers, detecting early prognostic biomarkers, predict-
ing their clinical significance, identifying mutational pat-
terns, and determining gene expression cohorts. The steps
in implementing the different techniques in integrating and
analyzing multi-omics data for cancer biomarker discovery
are illustrated in Fig. 2 [45]. Applying these different algo-
rithms for feature selection and classification and integrat-
ing high-dimensional heterogeneous omics data provides
potential inference methods in disease progressions, which
are listed in Table 1.
Supervised
Supervised ML algorithms are employed to assess labeled
datasets. These supervised methods analyze input and
output data across diverse classifications to train models.
These models, once trained, are subsequently utilized to
make predictions on multi-omics datasets. For instance,
they can be applied to identify the data characteristics that
underpin-specific biomarkers [46]. Supervised learning
algorithms are well-suited for addressing two fundamental
problems, which include classification and regression [42].
In the context of classification problems, the output vari-
able is distinctly discrete, necessitating the categorization
of diverse groups or categories of biomarkers with distinct
molecular features that are significantly altered between
healthy and diseased states [47]. In contrast, regression
problems involve output variables characterized by con-
tinuous real values, such as estimating the survival risk
in disease progression [48]. The versatility of supervised
learning algorithms positions them as valuable tools capa-
ble of effectively handling both discrete categorization and
continuous outcome prediction tasks.
Fig. 1   Overview of prevalent ML algorithms with attribute rankings [1—Low, 2—Medium, 3—High, 4—Very High] for concise representation
1273
Molecular Biotechnology (2025) 67:1269–1289
Logistic Regression
Logistic regression (LR) is a robust supervised classification
method, particularly applied in omics data analysis to iden-
tify and prioritize relevant biomarkers dynamically, optimiz-
ing the model’s predictive accuracy and interpretability in
the context of targeted therapies. Functioning as an exten-
sion of ordinary regression, LR models dichotomous varia-
bles, estimating the probability of an instance belonging to a
specific class ranging from 0 to 1 [49]. They are extensively
employed in assessing the associations between biomarkers
and binary outcomes and can be extended to handle the inte-
gration of diverse omics data types. Also, the LR algorithm
can seamlessly incorporate these diverse data types by treat-
ing them as covariates in the model, allowing the integra-
tion of multiple datasets. However, omics data frequently
exhibits high interdependence among features, violating the
assumption of independence in LR [50], leading to unstable
coefficient estimates and reduced interpretability. Research-
ers face the decision of categorizing biomarkers as either
continuous or categorical covariates. Predicting the survival
risk of the biomarkers with skewed distributions necessitates
normalization for integration as continuous covariates, often
accomplished through log transformation. Many biomark-
ers with skewed distributions align well with the Normal
curve on the log scale, recommending log transformation for
obtaining dependable odds ratio (OR) estimates predicting
clinical events [51]. Biomarkers identified through the omics
datasets encounter challenges in reproducibility owing to the
inherent heterogeneity stemming from diverse platforms or
laboratory settings [52]. These inconsistencies served as a
catalyst for the development of resilient biomarker identi-
fication methodologies through the integration of multiple
datasets. The distinct constant terms in LR gauged the sam-
ple heterogeneity. Through the minimization of variations in
constant terms within a given dataset, this approach effec-
tively preserved both intra-dataset homogeneity and inter-
dataset heterogeneity [53].
Support Vector Machines
Support vector machine (SVM) stands as a prevalently
used and robust ML algorithm in the field of bioinformat-
ics, particularly employed for classification and regression
tasks. The principle of SVM revolves around identifying
a hyperplane, either a line or a plane in high-dimensional
space that maximizes the separation between distinct classes
with the maximum margin. The crucial data points nearest
to this hyperplane, termed support vectors, wield significant
influence over its positioning. Researchers employ SVM on
Omics data for biomarker discovery, leveraging its sensi-
tivity and flexibility. The application involves metabolite
ranking through a systematic reduction of biological rep-
licates to assess the influence of sample size on biomarker
Fig. 2   Steps associated with implementing different strategies for integrating and analyzing multi-omics data for cancer biomarker discovery
1274
Molecular Biotechnology (2025) 67:1269–1289
reproducibility and robustness. This methodology reflects
SVM’s effectiveness in handling Omics data intricacies,
emphasizing its utility in the dynamic field of biomarker
identification and characterization [54]. SVM exhibits adapt-
ability to both linear and non-linear data, proving especially
efficacious when the number of features far surpasses the
number of samples. Despite its versatility, SVM presents
a spectrum of advantages and limitations, with its perfor-
mance intricately linked to the judicious choice of kernel
and the dataset’s size [55]. In multi-omics datasets, which
often involve intricate interactions [56], SVMs present an
advantageous method for detecting non-linear patterns that
might elude linear techniques like logistic regression, utiliz-
ing kernel functions. These kernel functions facilitate the
transformation of input data into higher-dimensional fea-
ture spaces, where the previously non-linear associations
between molecular attributes and clinical outcomes can
potentially become linearly separable. This transformation
enables SVMs to construct complex decision boundaries that
may not be achievable within the original input space. Com-
monly utilized kernel functions such as the radial basis func-
tion (RBF) kernel and polynomial kernel empower SVMs to
capture the intricate interactions among molecular features
[57]. Moreover, Ensemble methods, such as bagging and
boosting, can be integrated with SVMs to improve predic-
tive performance and robustness [58]. Consequently, SVMs
offer a robust avenue for uncovering concealed patterns and
relationships within multi-omics datasets, thus fostering
advancements in targeted therapy strategies.
Decision Tree
Decision tree (DT) stands as one of the earliest and most
prevalent ML algorithms, representing decision logic in a
tree-like structure to classify data. The tree comprises nodes
with multiple levels, starting with the root node at the top.
Internal nodes, conducting tests on input variables, guide
the classification process toward child nodes based on test
outcomes. This iterative process continues until it reaches
a leaf node, which signifies the decision outcome. Decision
trees are valued for their interpretability, quick learning, and
frequent integration into medical diagnostic protocols. Dur-
ing sample classification traversal, the outcomes of tests at
each node offer sufficient information to infer its class [59].
Table 1   The use of multiple ML approaches on OMICS datasets to identify novel biomarkers, early prognostic biomarkers, and prediction of
clinical significance
Omics data type
Assessment
Machine learning models
References
Transcriptomic
• Prognostic biomarkers identification
• Recurrence risk
• Intrinsic and extrinsic features of the disease environment
• Expression cohorts
• Autoencoder and Network
fusion
• Monte-Carlo feature
selection (MCFS)
• Support Vector Machine
(SVM)
• Principal component and
similarity measurements
[22, 102, 105–107]
Metabolomic
• Metabolic biomarkers in understanding disease pathophysiology
• Targeted therapeutic targets
• Pharmacogenomic responses
• Principle component
analysis and K-means
clustering
• SVM
• Random Forest
• KNN
• Naïve Bayes for classifier
models
[78, 108, 109]
Genomics
• Risk prediction
• Genetic mutation
• Mutational distribution
• Cancer somatic variants
• Copy number variation
• Understanding the genetic and molecular mechanisms
• Random Forest
• ANN
• Logistic regression
• SVM
• CNN
• DNN
[110–113]
Proteomics
• Survival risk prediction
• Protein levels, modifications, and interactions
• Biomarkers
• Linear models
• PCA
• K-Means clustering
[114, 115]
Epigenomics
• DNA methylation patterns
• Histone modifications
• Therapeutic response
• Alternative splicing prediction
• Disease classification
• PCA
• Naïve Bayes
• SVM
• Random forest
• RNN
[90, 116–118]
1275
Molecular Biotechnology (2025) 67:1269–1289
The DT undergoes initial training on “ground truth” data,
learning to establish decision boundaries for the most crucial
features based on class grouping. In the context of biomark-
ers, classes may represent cases and controls, with protein
levels serving as the features. Established methodologies
exist to identify the features contributing most significantly
to class differentiation. This not only enhances interpretabil-
ity but also mitigates the risk of classifying proteins that lack
medical relevance, avoiding potential confounding factors
related to distinct sample handling in cohort groups [60]. DT
can handle the non-linear relationships within omics data
analysis by employing recursive partitioning. This process
entails dividing the data into subsets according to predictor
variable values. At each node of the tree, decision trees dis-
cern the predictor variable and its corresponding split point
that can optimally segregate the data into homogeneous
groups relative to the outcome variable [61, 62]. Through
iterative data splitting based on predictor variables, decision
trees construct a hierarchical framework adept at modeling
non-linear interactions among molecular features in omics
datasets [63]. This adaptability renders decision trees well-
suited for delving into complex relationships and identifying
significant biomarkers or predictors within intricate biologi-
cal systems.
Random Forest
The random forest (RF) algorithm stands out as a premier
method for classification in ML due to its accuracy in han-
dling substantial datasets. It also operates as a learning algo-
rithm constructing multiple decision trees during training.
The collective predictions of individual trees contribute to
the model’s overall output. RF operates as a tree predictor,
with each tree relying on random vector values, showcasing
its efficacy in enhancing the robustness and predictive capa-
bilities crucial for accurate classification in bioinformatics
applications [64]. In omics analysis, employing systematic,
data-driven feature selection methods is crucial to avoid
biased selection. RF, in conjunction with other methods,
has proven effective in tasks such as gene and metabolite
selection for disease prognosis and progression [65–68].
RF inherently captures non-linear relationships between
molecular features and clinical outcomes, making it well-
suited for omics data analysis through partitioning the fea-
ture space into decision regions. RF can identify complex
patterns and interactions that may not be apparent with lin-
ear models, mitigating the risk of overfitting and noise often
encountered in high-dimensional omics datasets [69, 70].
This can be achieved by employing bagging, an ensemble
method, to enhance predictive accuracy and mitigate over-
fitting. This technique entails generating multiple bootstrap
samples from the original dataset and training multiple
decision trees on each sample. Through the integration of
predictions from numerous trees, RF diminishes variance
and achieves enhanced generalization to new data, thereby
averting overfitting [71]. Additionally, boosting, a sequen-
tial ensemble method, is employed, with subsequent models
learning from preceding misclassified models. Algorithms
like AdaBoost and Gradient Boosting Machines iteratively
refine the model, assigning greater importance to misclas-
sified instances, leading to improved overall performance
and diminished overfitting [71, 72]. The integration of these
ensemble methods with random forest can yield more resil-
ient models that exhibit superior generalization to high-
dimensional omics data while minimizing the susceptibility
to overfitting associated with specific training data patterns.
Existing studies often focus on specific omic types, lacking
stable feature selection procedures for power calculations
in identifying biomarkers. However, the lack of assessment
and validation of the identified markers hinders their utility
in study design or power analysis for translational research.
Naïve Bayes
The naïve bayes (NB) classifier, a probabilistic learning
model rooted in the Bayes theorem, is used for classifica-
tion tasks. This algorithm relies on the assumption of fea-
ture independence, making predictions about an instance’s
class by calculating the class prior probability and the likeli-
hood of that specific class. The NB classifier’s foundation
in probability theory renders it valuable for diverse clas-
sification bioinformatics applications [73]. The algorithm
can be described using: P(X|Y) = P(Y|X)P(X)∕P(Y) . The
P(X|Y) represents the Posterior, indicating the probability of
X being true given that Y is true. On the other hand, P(Y|X)
describes the likelihood of a class, representing the prob-
ability of B being true given that A is true. Additionally,
P(X) and P(Y) denote the prior probability of a class and a
predictor, respectively. NB algorithms find application in
the selection of genetic biomarkers through the concurrent
examination of genome-wide SNP data and large omics data
[74, 75]. NB algorithm has not been extensively explored
for predicting biological classes. However, refined Bayesian
classification methods that consider dependencies among
features have demonstrated precise predictions of biological
classes. NB effectively handles noisy and irrelevant features
in the high-dimensional data by leveraging its probabilistic
framework. Noisy features have minimal impact on condi-
tional probabilities, as NB considers joint probability distri-
butions. Similarly, irrelevant features insignificantly affect
class probability estimation, as NB prioritizes discriminating
between classes based on informative features in the omics
data [76]. Nevertheless, enhancing the robustness of predic-
tions is achieved by employing an ensemble approach, spe-
cifically bagging, with Naïve Bayes classifiers. This strategy
enhances the effectiveness of ranking and selecting attributes
1276
Molecular Biotechnology (2025) 67:1269–1289
used by each bagged classifier, ultimately reinforcing attrib-
ute independence in the biomarker selection process [77].
K‑Nearest Neighbours
K-nearest neighbors (KNN) is a versatile and efficient ML
algorithm suitable for classification and regression tasks. It
classifies an unknown sample based on its proximity to the
K-nearest samples in the training set, assigning the most
common class among these neighbors. As a lazy learning
algorithm, KNN stores training data during the training pro-
cess, conducting actual classification or regression at the
prediction stage for enhanced speed and memory efficiency.
Despite its simplicity, adaptability to linear and non-linear
data, and ease of implementation, KNN’s performance
depends on parameters like the selection of K, feature scale,
and the relevance of features [55]. The KNN algorithm is
applied to discern patterns within high-dimensional omics
datasets, classify or predict sample phenotypes based on
proximity in feature space, and aid in identifying potential
biomarkers by revealing similarities among biological sam-
ples. Researchers exploit KNN’s adaptability and simplicity
to explore intricate omics relationships, contributing to bio-
molecular marker discovery in precision medicine [78–80].
KNN’s non-parametric nature allows it to capture complex
relationships and patterns within multi-omics datasets with-
out making strong assumptions about data distribution [81].
This attribute is particularly valuable in deciphering intricate
molecular landscapes associated with disease phenotypes,
areas where conventional linear models may encounter limi-
tations. KNN is frequently employed for imputing missing
metabolite abundances in omics datasets [82]. Nevertheless,
it is crucial to acknowledge that KNN operates under the
assumption that missing values are uniformly distributed at
random across the dataset, a premise that does not align
with the typical characteristics of metabolomics data [83].
Despite this consideration, the algorithm’s versatility in
handling classification and regression tasks integrated with
Artificial Intelligence (AI) methodologies can enhance its
utility for uncovering novel insights in personalized health-
care [84].
Artificial Neural Network
Artificial neural networks (ANNs) draw inspiration from
biological neural networks, resembling interconnected artifi-
cial neurons. These artificial neurons receive input, undergo
a data transformation, and produce an output, mirroring the
functional principles of their biological counterparts. The
ANN model does not make any assumptions about the dis-
tribution of data before the learning process, enhancing the
versatility and applicability of ANNs across various domains
[85]. ANNs comprise input and output layers connected by
hidden layers. Input nodes transmit information to hidden
nodes through activation functions, while hidden nodes acti-
vate based on presented evidence. Weighting functions in
hidden layers process evidence, and when node values reach
a threshold, outputs are generated. ANNs require extensive
training data, limiting application in rare events with insuf-
ficient data. They do not accommodate human expertise
substitution for quantitative evidence. The key advantages
of employing ANNs are that they exhibit robust fault and
failure tolerance, scalability, and reliable generalization
capacity, enabling accurate prediction or classification of
novel, ambiguous, and unlearned data, which makes it suit-
able for biomarker studies, contributing to the development
of biomarker panels that, when used collectively, enhance
prognostic capabilities in diseases. ANNs trained on large-
scale multi-omics datasets can be adapted and transferred to
new domains or disease contexts with limited labeled data.
Transfer learning techniques enable the transfer of knowl-
edge learned from one dataset to another, accelerating model
training and improving the predictive performance of ANNs
in multi-omics data analysis [86]. However, ANNs can be
integrated with dimensionality reduction methods like
autoencoders or t-distributed stochastic neighbor embed-
ding (t-SNE) to reduce the dimensions of omics data while
retaining critical features. This enables the visualization of
high-dimensional data in lower-dimensional spaces, which
assists in deciphering intricate molecular relationships [87].
Deep Neural Networks
Deep neural networks (DNNs) serve as potent tools for
data-driven modeling in bioinformatics. Comprising layers
with interconnected nodes and edges that encapsulate math-
ematical relationships, DNNs undergo iterative refinement
via backpropagation during training. Post-training, these
updated relationships function as predictive equations, ena-
bling the accurate forecasting of output variables based on
input variables. An inherent strength of DNNs lies in their
ability to capture and express intricate relationships within
a system, irrespective of its non-linear and complex nature
[88]. The intricate and interlinked hierarchical representa-
tions of training data utilized by deep neural networks for
estimating purposes render the comprehension of these esti-
mates exceptionally challenging, resulting in low explain-
ability. However, we identified two deep learning neural net-
works, Convolution Neural Networks (CNNs) and Recurrent
Neural Networks (RNNs) that are widely used for Omics
integration and analysis. The CNNs are found to recognize
spatial patterns in sequences and are efficient for identify-
ing relevant genetic markers [89]. Moreover, RNNs capture
temporal dependencies in time-series omics data, aiding
in discerning dynamic biomarker patterns [90]. The sig-
nificance of DNNs in omics data analysis is underscored by
1277
Molecular Biotechnology (2025) 67:1269–1289
the burgeoning advancements in multi-omics technologies
and the accumulation of extensive bio-datasets with issues
like overfitting, interpretability deficits, data heterogeneity
integration challenges, and the need for enhanced predic-
tion accuracy. Integrated approaches with dimensionality
reduction methods are established to extract targets from
each omics data and construct sample similarity networks
based on feature matrices. Subsequently, the fused similarity
networks undergo training in a DNN, significantly reducing
data dimensionality and mitigating the risk of overfitting
[91]. This advancement holds promise for the evolution of
targeted therapy.
Unsupervised
Unsupervised ML is employed for the data lacking labels.
The model uncovers latent patterns by identifying groups
of samples sharing common characteristics. The speed and
reliability of the unsupervised ML methods depend on the
nature of the data and are sensitive to outliers. Method selec-
tion balances computational efficiency and captures specific
patterns in complex data. However, Accuracy is gauged by
the likelihood of the model generating a dataset under a
specified distribution. Unsupervised learning aims to clus-
ter data, reveal natural groupings, or establish associations
among data points within large databases [92]. Unsuper-
vised multi-omics methods primarily focus on classifying
diseases and sample subtypes while identifying biomarkers
or modules associated with a disease. Current methods often
handle multiple outcome variables individually instead of
using multivariate models.
Principle Component Analysis
Dimension reduction is vital for downstream tasks like
pattern recognition, classification, and clustering in high-
dimensional data [55]. While various techniques exist, prin-
cipal component analysis (PCA) stands out as a classical
and widely employed method. PCA offers optimal linear
projection in Euclidean space with eigenvectors represent-
ing weighted linear combinations of features [93]. PCA
can function as a valuable asset for quality control and the
mitigation of batch effects in multi-omics investigations.
Through the visualization of data distributions and the detec-
tion of outliers or batch effects, PCA empowers research-
ers to evaluate data quality and ensure consistency across
various experimental batches or sample groups [94]. Con-
sequently, this process contributes to bolstering the reliabil-
ity and reproducibility of downstream analyses. Identifying
principal components associated with disease phenotypes
can prioritize molecular features that contribute most sig-
nificantly to the observed variation, guiding the selection
of potential targets. Acknowledging that only a subset of
features is dominantly relevant in genomic transcriptomic
studies, using all features may compromise robustness and
interpretability in high-dimensional data. Therefore, Park
et al. introduced an integrative analysis of multi-omics data,
utilizing blockwise sparse principal components to alleviate
multi-collinearity and redundancy, achieve dimensionality
reduction, and identify crucial variables, unlike other multi-
omics integration methods for biomarker discovery [95].
K‑Means Clustering
The K-means clustering divides a data space into k clus-
ters, assigning each data point to the cluster with the nearest
mean value. The clustering involves partitioning data into
groups with similar characteristics, particularly focusing
on geometric proximity in the feature space. This approach
aids in creating a learning problem for precise recovery of
cluster centroids, eliminating impractical considerations
[96]. K-means clustering is sensitive to the initial selection
of cluster centroids, which can lead to suboptimal solutions
and influence cluster assignments [97]. Additionally, it oper-
ates under the assumption of equal variance among clusters,
which may not hold for multi-omics data with varying levels
of noise and biological variability. To address these limita-
tions, performing multiple initializations of K-means with
different seed values and opting for the solution exhibiting
the lowest within-cluster variance can mitigate the sensitiv-
ity to initialization, enhancing the robustness of clustering
outcomes. Integrating semi-supervised learning techniques
with K-means clustering can leverage labeled data and prior
knowledge in multi-omics analysis, improving the interpret-
ability and accuracy of subtype identification and biomarker
discovery [98]. Moreover, recent research has delved into
the analysis of multilayer data and similar studies, demon-
strating enhanced predictive capabilities for disease outcome
models compared to single-layer analyses. K-Means primar-
ily segregate distinct data types but struggle to group various
interconnected omics measurements into cohesive clusters
[99, 100].
Autoencoders
The autoencoder functions as a neural network that aims
to replicate the input signal, capturing essential features of
the data and restoring its original form. It employs a hidden
layer as both an encoder and decoder, ensuring consistency
between the encoded and decoded data. Autoencoders often
utilize greedy layer-wise pretraining for unsupervised learn-
ing. Commonly employed in scenarios with limited labeled
data, autoencoders serve omics data exploitation well, given
the challenges of obtaining labeled omics data, typically
characterized by high dimensionality [101, 102]. Adversar-
ial training generates adversarial examples that encourage
1278
Molecular Biotechnology (2025) 67:1269–1289
the autoencoder models to acquire more discriminative and
stable representations of multi-omics data, particularly in
the presence of noise and variability [103]. However, inter-
preting the latent features learned by autoencoders can be
challenging, as they often represent abstract combinations
of molecular attributes. This lack of interpretability may
hinder the biological insights gained from clustering analy-
sis and limit the utility of autoencoders on omics datasets.
Therefore, implementing regularization techniques such as
dropout and weight decay can mitigate overfitting and pre-
vent the model from learning noise and irrelevant patterns
in the data, which can enhance the robustness of clustering
outcomes [104].
Multi‑omics Data Visualization and Analysis
Drug target identification and prioritization are critical
challenges in the pre-clinical stages of pharmaceutical
research. Computational techniques can reap the benefits
of relatively large human genomics and proteomics data
to identify targets, minimizing the time and cost. The
paradigm of the “one drug target–one disease” approach
in drug discovery has become inefficient. The emergence
of phenotypic resilience and network topology by break-
throughs in systems biology strongly suggests that pre-
cisely selective molecules may have lower clinical efficacy
than multi-target treatments. The network pharmacology
approach has lately acquired prominence as a strategy for
integrating omics data and developing multi-target drugs,
respectively [119]. Networks can be defined at many lev-
els of complexity (Fig. 3). Protein–protein interaction,
gene regulatory, and metabolic networks are common
examples of biological networks. The network concept is
sometimes expanded to include drug–drug interaction in
modern pharmacological research. Diverse forms of data
will lead to distinct network features in terms of linkage,
complexity, and structure, with edges and nodes possibly
conveying many layers of data. Polypharmacology is more
appropriate for complex diseases that involve complex
target networks and biological pathways [120]. Network
pharmacology finds diverse applications in addressing
research questions. It serves various purposes, such as
uncovering disease-causing candidate genes, revealing dis-
ease-associated subnetworks and systematic perturbations,
and capturing therapeutic responses for efficient target
identification and drug discovery. Integrating omics data
into static and dynamic network models in cancer enables
the identification of key molecular players, dysregulated
pathways, and potential therapeutic targets. These models
can also predict drug responses, guide personalized treat-
ment strategies, and explore novel therapeutic interven-
tions. Table 2 lists different cancer network models with
their applications in integrating omics data.
Fig. 3   Classification of network models integrating multi-omics data and depicting the complexity of the models
1279
Molecular Biotechnology (2025) 67:1269–1289
Table 2   The applications and tools employed for different network models with integration of omics data for prognostic biomarkers and therapeutic target identification
Network model
Omics data integration
Applications
Tools
References
Protein–Protein Interaction Networks
(PPINs)
Integration of gene expression, protein
abundance, or proteomics data with
PPINs
Identifying key proteins, functional mod-
ules, and signaling pathways associated.
Discovery of potential therapeutic targets
and biomarkers
• STRING (Search Tool for the Retrieval
of Interacting Genes/Proteins)
• Dynamic Protein–Protein Network with
Support Vector Machine (DPPN-SVM)
(https://​github.​com/​brown-2/​mis\_​local​
izati​on)
[142, 143]
Gene Regulatory Networks (GRNs)
Integration of gene expression, transcrip-
tion factor binding, and miRNA expres-
sion data with GRNs
Understanding gene regulation mecha-
nisms. Identification of dysregulated
pathways and regulatory modules
NetworkAnalyst (https://​www.​netwo​rkana​
lyst.​ca/)
[144, 145]
Signaling Pathway Networks
Integration of phosphoproteomics, tran-
scriptomics, or proteomics data with
signaling pathway networks
Investigation of dysregulated signaling
pathways. Identification of potential
therapeutic targets
CNORode—a logic-based ordinary differ-
ential equation (https://​github.​com/​saezl​
ab/​CNORo​de)
[135, 146]
Metabolic Networks
Integration of metabolomics, fluxomics,
or transcriptomics data with metabolic
networks
Understanding altered metabolic pathways.
Identification of metabolic rewiring and
potential targets for intervention
• metabolic Context-specificity Assessed
by Deterministic Reaction Evaluation
(mCADRE) (https://​github.​com/​jaeddy/​
mcadre)
• Reconstruction, Analysis, and Visualiza-
tion of Metabolic Networks (RAVEN)
[147–150]
Co-expression Networks
Integration of gene expression or transcrip-
tomics data with co-expression networks
Identification of functionally related genes.
Discovery of co-expression patterns and
potential biomarkers
• Weighted Gene Co-Expression Network
Analysis (WCGNA)
• CoExp
[151–154]
miRNA-gene Regulatory Networks
Integration of miRNA expression and gene
expression data with regulatory networks
Investigation of miRNA-mediated regula-
tory interactions. Identification of dys-
regulated miRNAs and their target genes
NetworkAnalyst (https://​www.​netwo​rkana​
lyst.​ca/)
[145, 155]
Gene Co-mutation Networks
Integration of mutation data (e.g., somatic
mutation, copy number variations) with
co-mutation networks
Identification of co-occurring genetic muta-
tions and understanding the landscape of
genetic alterations
Pajek (http://​vlado.​fmf.​uni-​lj.​si/​pub/​netwo​
rks/​pajek/)
[156, 157]
Drug–gene Interaction Networks
Integration of drug response data, gene
expression, or genomic data with drug–
gene interaction networks
Prediction of drug response and identifica-
tion of potential drug targets in cancer
treatment
Drug–Gene Interaction Database (DGIdb)
(https://​www.​dgidb.​org/)
[158, 159]
1280
Molecular Biotechnology (2025) 67:1269–1289
Static Networks
The advent of high-throughput omics technologies has
provided abundant data that can be integrated into static
network models to enhance our understanding of complex
diseases. Omics data from genomics, such as somatic
mutations and copy number alterations, can be integrated
into static network models to identify driver genes and
elucidate disrupted signaling pathways [121–123]. Gene
expression profiles derived from transcriptomics data
enable the construction of gene co-expression networks,
unraveling functional modules and identifying poten-
tial biomarkers associated with specific cancer subtypes
or stages [124–126]. By incorporating proteomics data,
static network models can capture protein–protein interac-
tions, shedding light on key protein hubs, signaling cas-
cades, and potential therapeutic targets [127]. Integrating
metabolomics data into metabolic network models allows
for identifying altered metabolic pathways, facilitating
the discovery of metabolic vulnerabilities and potential
avenues for therapeutic intervention [128]. Computational
approaches, including correlation-based algorithms and
mutual information analysis, enable the inference of regu-
latory relationships and interactions between molecular
components based on omics data [129]. Advanced visuali-
zation techniques, such as network graphs, aid in compre-
hending the intricate structure of static network models.
Network analysis methods, such as centrality measures and
module detection algorithms, assist in identifying criti-
cal nodes, dysregulated pathways, and functional modules
within the network [130, 131]. A set of molecular regula-
tors (genes or parts of genes) that interact with one another
and other components in the cell to control the levels of
gene expression of mRNA and proteins, which define the
cell’s function, is called gene regulatory networks. GRNs
have significant challenges, such as the delay between
transcription and translation and a lack of knowledge of
the necessary topology to depict the targeted phenotype
and kinetics [132].
Static network models provide a comprehensive systems
level view of cancer’s intricate molecular interactions and
regulatory relationships, integrating multiple layers of omics
data. This facilitates the generation of testable hypotheses
concerning the functional roles of genes, proteins, and path-
ways, thereby aiding in the discovery of novel therapeutic
targets and biomarkers. Static network models lack tempo-
ral information and fail to capture the dynamic changes in
molecular interactions over time, impeding a comprehen-
sive understanding of disease progression and treatment
responses. They possess limitations regarding coverage of
the entire molecular landscape, potentially overlooking criti-
cal genes, proteins, or interactions. They often oversimplify
the heterogeneity and contextual variability, treating it as
a homogeneous entity and neglecting the diversity across
disease subtypes or individual patients.
Dynamic Networks
Dynamic network models have emerged as powerful tools
for understanding temporal and disease-specific dynamics.
By integrating multi-dimensional omics data into these mod-
els, researchers can capture the intricate regulatory events,
signaling pathways, and molecular interactions that drive
disease progression and therapeutic responses. Differential
equations and logic functions are generally used to gener-
ate dynamic systems [133, 134]. Dynamic models assist in
evaluating the cumulative effect of dysregulated molecular
mechanisms in an individual, guiding therapeutic choice in a
targeted approach. Simulating metabolic and regulatory net-
works usually entails non-linear and dynamic models [135].
However, due to technical constraints, such models typically
focus on one or a few metabolic pathways: for example, a
dearth of reaction kinetics and mechanism data, the high
computing complexity of non-linear parameter estimation,
and a scarcity of dynamic experimental data for simulation
validation [136].
On the other hand, Stoichiometric models are usually
genome-wide and generated from metabolic networks,
allowing them to be easily combined with high-throughput
data. Genome-scale metabolic models (GEMs) are concep-
tual mathematical models that assist the investigators in
estimating the metabolic flux rates to determine metabolic
alterations to meet increased energy needs for growth, sur-
vival, rapid proliferation, and other features of tumor cells
in various cancer pathogenesis through simulation studies
and hypothesis testing [137, 138]. Also, cancer cells can be
precisely targeted using genetically engineered therapeutics
that take control of specific driving pathways through gene
interaction networks (GINs) that emphasize an unbiased
assessment of dysregulated pathways and uncover genetic
variation that may be used to design targeted therapeutics
[139].
Molecular assessment at the systems level requires cap-
turing the static protein activity in the cell, with all interac-
tions uncovered being integrated into a single data structure
to derive protein–protein interaction networks (PPINs).
Recent systems biology trends aim to develop tailored
PPINs representing specific conditions [140]. Zhang et al.
generated dynamic PPINs by integrating high-throughput
and gene expression data to determine certain proteins’
active probability and time points. Dynamic PPINs, in con-
trast to static PPINs, may efficiently express both dynamic,
active, and topological information in PPINs [141]. The
dynamic models require time-resolved or context-specific
omics data, which can be challenging to obtain and subject
to noise and experimental limitations. It requires parameter
1281
Molecular Biotechnology (2025) 67:1269–1289
estimation, such as the initial conditions, reaction, or decay
rates [128]. Estimating these parameters accurately from
limited experimental data can introduce uncertainty into the
model. Inaccurate or uncertain parameter values can affect
the model’s predictive power and limit its ability to capture
the true dynamics of the biological system. Addressing these
flaws and limitations requires continuous improvement in
data collection techniques, computational algorithms, and
integration of multiple omics data types. Efforts to incor-
porate more comprehensive and context-specific data, con-
sider non-linear relationships and account for variability
across diseases and patients can enhance the accuracy of
the dynamic network models.
Integrative Approach for Probing Disease
Mechanisms
Biological processes and molecular functions arise from
intricate interactions among thousands of molecules, con-
stituting inherent complexity. Integration of metabolomics
data with other omics data holds significant promise for
achieving a holistic understanding of disease mechanisms.
Metabolomics, which focuses on the comprehensive analysis
of small molecule metabolites within biological systems,
provides unique insights into the functional status and meta-
bolic phenotypes associated with various physiological and
pathological conditions [160, 161]. The integration of omics
datasets with computational models and network analysis
tools elucidates the complex interplay between genes, pro-
teins, metabolites, and cellular processes underlying disease
phenotypes.
Despite recent progress in omics technologies, the under-
lying genetic factors contributing to numerous metabolic
phenotypes remain elusive. Metabolite biomarkers can be
integrated with genomics and clinical parameters to enhance
diagnostic accuracy or refine disease risk prediction mod-
els. Metabolites can also serve as intermediate phenotypes
for genetic investigations, offering insights into underlying
genetic mechanisms [162]. The integration of metabolomics
data with either whole-exome sequencing or WGS-data pre-
sents a promising systematic strategy for pinpointing dis-
ease-causing variants and holds potential utility within the
framework of a specific pathway under investigation [163].
Furthermore, at a more intricate biological and analytical
level, metabolomics can be combined with various omic
platforms, facilitating a comprehensive understanding of
complex biological systems and interactions (Fig. 4).
The alterations in metabolite levels, perturbations in met-
abolic pathways, and the onset of disease states can be eluci-
dated by assessing the epigenetic alterations. This approach
offers molecular insights into the intricate interplay among
genetic, epigenetic, and metabolic factors during the dis-
ease progression. Through the integration of epigenomic
Fig. 4   The workflow for integration of metabolomics with other omics for a holistic understanding of disease progression
1282
Molecular Biotechnology (2025) 67:1269–1289
and metabolomic data, the intricate relationships between
epigenetic alterations and metabolic pathways in disease
pathogenesis can be uncovered. In recent years, metabo-
lomics and epigenomics have experienced notable advance-
ment as prominent molecular and analytical methodologies
for biomarker identification [164, 165]. In the context of
cancer, it is characterized by distinctive features such as
metabolic reprogramming and epigenetic modifications,
which play pivotal roles in tumor progression and are intri-
cately interconnected with the tumor microenvironment and
other molecular pathways [166]. Epigenetic modifications
can directly influence the expression of metabolic genes,
thereby altering cellular metabolism and contributing to dis-
ease phenotypes [167]. Conversely, metabolic alterations can
impact epigenetic regulation by modulating the availability
of metabolites involved in epigenetic modifications [166,
168]. The cross-talk between epigenetics and metabolism
represents a dynamic interplay that shapes cellular physiol-
ogy and disease susceptibility [169, 170]. DNA methyla-
tion stands out as a widely studied epigenetic mechanism
with implications for cancer-related gene regulation. Meth-
ylation processes, occurring directly in promoter regions of
cancer-related genes or histone residues, significantly influ-
ence DNA accessibility and gene expression regulation. The
availability of methyl groups, primarily mediated by metabo-
lites within the methionine and folate cycles, closely relates
to DNA methylation processes [171]. Alterations in meta-
bolic concentrations of the Tricarboxylic Acid (TCA) cycle
intermediates, such as α-ketoglutarate (α-KG), succinate,
fumarate, and acetyl-CoA, significantly impact chromatin-
modifying enzymes, including 10–11 translocation enzymes
and histone demethylases. These enzymes play crucial roles
in catalyzing hydroxylation and demethylation processes,
ultimately shaping the epigenetic landscapes of cancer.
Additionally, metabolites such as succinate, fumarate, and
α-KG, termed oncometabolites, accumulate in cancer cells
due to mutations in fumarate hydratase and succinate dehy-
drogenase genes, further influencing epigenetic alterations
[170, 172]. Overall, the regulation of tumor gene expression
reflects a complex interplay between epigenetic enzymes and
metabolic substrates, demonstrating the intricate mecha-
nisms underlying cancer pathogenesis.
The conventional linear model of data progression from
genes to transcripts, proteins, and metabolites is being recon-
sidered to recognize the intricate association of these net-
work layers. Relying solely on single-level data often proves
inadequate for fully understanding biological processes. For
instance, fluctuations in metabolite concentrations may arise
from downstream production or reduced metabolism from
upstream reactions, making causal assessments challeng-
ing. Similarly, while transcripts provide insight into gene
transcription, they do not convey the functionality or activ-
ity of resulting proteins [173]. Consequently, integrated
Omics methodologies with transcriptomics and metabo-
lomics emerging as commonly preferred combinations in
contemporary investigations are being adopted to gain a
more comprehensive understanding of disease progression
with novel regulatory pathways and biomarkers [174]. This
approach yields comprehensive datasets elucidating the
interconnected metabolic and transcriptomic alterations.
This integration will facilitate the identification of relation-
ships between proteins and metabolites, thereby revealing
molecular me
\end{document}
